% Sections on the nine acceleration proposals, fastunknot 0.2 and the Rust port.
% Included from report.tex.  Tables are generated by make_tables.py.

\section{Nine acceleration proposals}
\label{sec:proposals}

\subsection{What was delivered}
After version 0.1 of \code{fastunknot} (Section \ref{sec:fast}) was written, nine
further runs were each asked to accelerate it. Their archives were extracted
into \code{proposals/01}--\code{09} (wrapper directories removed, ZIP files
deleted, archive order recorded in \code{proposals/ORDER.txt}). Every proposal
is a complete replacement package with its own report, tests, an unchanged
copy of the 0.1 baseline and same-machine benchmarks, and every one of them
states that the general worst case remains exponential. As with the six
original archives, the independent runs converged on the same ideas, which is
evidence that the ideas are the natural ones.

\subsection{The ideas}
\begin{enumerate}
\item \textbf{Units are involutions.} In
  $\operatorname{End}(m)\cong\F[x_1,\dots,x_k]/(x_i^2)$ a unit is $1+\nu$ with
  $\nu$ in the maximal ideal, and $\nu^2=0$ in characteristic two (squares of
  monomials vanish, cross terms occur twice). So $(1+\nu)^{-1}=1+\nu$ and the
  geometric series of 0.1 is unnecessary. All nine.
\item \textbf{Min-fill pivots.} Version 0.1 cancelled invertible entries in
  last-in-first-out order. Choosing the pivot $b\to c$ that minimizes the
  Markowitz count $(\mathrm{out}(b)-1)(\mathrm{in}(c)-1)$ through a lazy heap
  limits fill-in \cite{markowitz}. Seven of nine (not 05, 06).
\item \textbf{Cheaper cobordism algebra.} The topology of a gluing (which discs
  merge, Euler characteristics, boundary circles) does not depend on the dots,
  so it can be computed once per triple of matchings and reused; identity and
  square shortcuts; memoized results. All nine; proposal 03 goes furthest and
  packs a morphism into one integer whose bit number is a dot mask.
\item \textbf{$O(n\log n)$ scan order and $O(n)$ descending test.} The greedy
  score of a crossing changes only when a neighbour is processed, so a lazy
  heap reproduces the 0.1 order exactly; the descending test becomes a sweep
  over a difference array of bad cyclic intervals. All nine.
\item \textbf{Modular Alexander test.} Evaluate the first Alexander minor in a
  prime field; for the unknot it is $\pm t^k$, $0\le k<n$, so any other value
  certifies knottedness in $O(n^3)$ field operations. Six of nine, three of
  them only at $t=-1$ (the determinant).
\item \textbf{Modular Jones test by a frontier scan.} Evaluate the Kauffman
  bracket at a fixed $A\in\mathbb F_p$ while scanning the diagram, adding the
  weights of partial states as soon as they induce the same boundary matching.
  A value different from $\delta(-A^3)^{w}$ certifies knottedness. All nine.
  This is the most valuable single idea for recognition: it rejects the
  Kinoshita--Terasaka and Conway knots, which have trivial Alexander
  polynomial, in about a millisecond, and its cost is \emph{provably}
  $\mathrm{Catalan}(w/2)\cdot\mathrm{poly}(n)$ for scan width $w$, because
  states with equal matchings are merged; there is no multiplicity factor.
\item \textbf{Visible connected sums.} Two edges bordering the same two faces
  are a two-edge cut of the projection, so the diagram is a connected sum of
  the two sides. Reduced Khovanov homology over a field is multiplicative, so
  ranks multiply and degree counts convolve; for recognition the summands are
  examined separately. Seven of nine.
\item \textbf{Crossing signs from both strands.} Version 0.1 assumed that a PD
  row starts at the incoming under-port (true for its own converters); seven
  proposals compute the sign from both incoming ports, which the Jones
  normalization needs.
\item \textbf{Incremental Reidemeister simplification} (02 only), \textbf{no
  cancellation on the last crossings followed by linear algebra} (06 only),
  \textbf{an exact integer bracket after the modular one} (05 only), and
  \textbf{witness replay commands} (six proposals).
\end{enumerate}

\subsection{Same-machine comparison}
The harness \code{synthesis/data/proposals\_bench.py} runs every engine on the
same inputs, each measurement in a fresh process with a 120-second limit;
timers exclude start-up, imports and input construction; medians of three runs
(one run above five seconds). ``scan'' is the raw Khovanov scan with
factorization disabled, ``rank'' the rank of a connected sum of eight Conway
knots through each engine's factored API when it has one, ``recognize'' the
default pipeline. Milliseconds; 0.1 is the baseline, 0.2 the integrated
version of Section \ref{sec:v02}, and Rust the port of Section \ref{sec:rust}.

\begin{center}\scriptsize\setlength{\tabcolsep}{3.2pt}
\IfFileExists{tables/proposals_bench.tex}{\begin{tabular}{@{}lrrrrrrrrrrrr@{}}
\toprule
Task & 0.1 & 01 & 02 & 03 & 04 & 05 & 06 & 07 & 08 & 09 & 0.2 & Rust \\
\midrule
scan: conway & 34.4 & 11.6 & 10.4 & 9.86 & 11.3 & 10.2 & 10.4 & 9.80 & 10.0 & 12.5 & 8.44 & -- \\
scan: kinoshita\_terasaka & 41.5 & 12.4 & 11.2 & 10.5 & 12.8 & 12.3 & 13.3 & 10.8 & 11.3 & 13.0 & 9.05 & -- \\
scan: T(3,11) & 51.8 & 29.9 & 29.4 & 27.9 & 31.2 & 28.8 & 31.7 & 26.9 & 30.3 & 30.8 & 26.8 & -- \\
scan: braid3\_40 & 611 & 103 & 115 & 94.0 & 108 & 133 & 140 & 143 & 164 & 201 & 84.1 & -- \\
scan: braid4\_41 & 984 & 263 & 274 & 212 & 350 & 520 & 321 & 273 & 268 & 284 & 160 & -- \\
scan: braid6\_31 & 495 & 96.7 & 108 & 88.2 & 115 & 95.1 & 93.6 & 94.6 & 104 & 114 & 67.7 & -- \\
scan: braid5\_36 & $>$120\,s & 1064 & 1208 & 887 & 1153 & 44995 & 7708 & 1080 & 1163 & 1491 & 691 & -- \\
scan: chain256 & 312 & 28.8 & 24.9 & 29.7 & 26.9 & 24.8 & 335 & 24.1 & 24.2 & 25.3 & 19.3 & -- \\
scan: conway\_sum\_2 & 1481 & 239 & 241 & 173 & 204 & 220 & 205 & 213 & 207 & 261 & 136 & -- \\
\midrule
rank: conway\_sum\_8 & $>$120\,s & 31.5 & $>$120\,s & 20.2 & $>$120\,s & $>$120\,s & 16.0 & $>$120\,s & 51.7 & 18.6 & $>$120\,s & -- \\
\midrule
recognize: trefoil & 0.13 & 0.19 & 0.19 & 0.13 & 1.67 & 0.15 & 0.15 & 0.28 & 0.16 & 0.15 & 0.27 & -- \\
recognize: torus\_3\_5 & 0.98 & 0.34 & 1.03 & 0.32 & 2.24 & 1.06 & 0.80 & 1.04 & 0.32 & 0.66 & 0.56 & -- \\
recognize: hard\_unknot\_8 & 2.84 & 5.05 & 2.40 & 4.02 & 3.86 & 2.45 & 3.07 & 2.18 & 2.16 & 3.03 & 2.78 & -- \\
recognize: unknot\_braid40 & 11.2 & 8.00 & 0.41 & 7.82 & 8.31 & 7.64 & 7.50 & 7.93 & 8.64 & 8.42 & 0.17 & -- \\
recognize: conway & 58.8 & 1.53 & 0.89 & 1.79 & 4.01 & 3.50 & 4.05 & 1.57 & 3.95 & 1.20 & 1.27 & -- \\
recognize: kinoshita\_terasaka & 65.6 & 2.12 & 1.02 & 1.49 & 3.43 & 1.55 & 3.31 & 1.08 & 1.81 & 0.72 & 1.25 & -- \\
recognize: conway\_sum\_2 & 1817 & 3.54 & 2.05 & 3.47 & 2.75 & 2.32 & 12.6 & 2.07 & 2.77 & 2.17 & 1.56 & -- \\
recognize: conway\_sum\_3 & 58090 & 6.77 & 2.96 & 5.15 & 3.60 & 3.43 & 35.3 & 2.10 & 2.78 & 2.51 & 2.47 & -- \\
recognize: T(3,61) & 4747 & 3.71 & 12.5 & 5780 & 13.8 & 19.0 & 5811 & 22.3 & 5.91 & 8.80 & 8.84 & -- \\
recognize: braid5\_36 & 9.89 & 2.53 & 0.69 & 2.38 & 8.63 & 2.63 & 11.2 & 6.20 & 2.40 & 2.37 & 0.90 & -- \\
\bottomrule
\end{tabular}
}{(table pending)}
\end{center}

All eleven Python engines and the Rust binary return identical ranks and
verdicts on every input. Observations:
\begin{itemize}
\item The stress case \code{braid5\_36}, which 0.1 abandoned after 600 seconds,
  takes about a second wherever min-fill pivots are used, 45 seconds in
  proposal 05 (LIFO kept) and 7.7 seconds in proposal 06 (LIFO plus the tail
  trick). Its reduced rank is 2949. \textbf{This corrects our own 0.1
  analysis}: we had attributed the timeout to large minimal complexes of the
  intermediate tangles. The minimal complexes are small (at most 5898
  objects); what exploded was fill-in caused by the pivot order.
\item Proposal 03 has the fastest raw scanner, which is why its
  representation was adopted.
\item On $T(3,61)$ proposals 03 and 06 take 5.8 seconds. Proposal 06 runs the
  symbolic Alexander polynomial before anything modular. Proposal 03 is more
  instructive: it evaluates the minor at $t=2$ modulo $2^{31}-1$, where $2$
  has multiplicative order 31, so the set of ``units'' $\{\pm2^k\}$ has only
  62 elements and is closed under the cyclotomic arithmetic of torus knots;
  $\Delta_{T(3,61)}(2)\equiv-2^{29}$ and the filter misses. Our own first
  implementation had the same structural weakness ($2$ has order 61 modulo
  $2^{61}-1$) even though it happened to fire. Version 0.2 and the Rust port
  evaluate at fixed generic field elements instead, for both the Alexander and
  the Jones test. The filters are one-sided, so this was never a soundness
  issue, only a lost speedup.
\item Tiny inputs get slower: the trefoil goes from 0.1 ms to 0.2--0.3 ms in
  most engines, because more stages run before one fires.
\end{itemize}

\section{\code{fastunknot} 0.2}
\label{sec:v02}

\subsection{What was adopted}
Version 0.2 (\code{fast/}) integrates every proposed idea that produced a
measurable speedup, re-implemented in the structure of the existing package;
the 0.1 scanner is kept unchanged as \code{scan\_reference} so that each idea
can be switched off and so that the new code has an independent oracle.

The pipeline is now: validation; incremental R1/R2 reduction; linear descending
test; visible connected-sum factorization (summands examined smallest first,
the first knotted one decides, and a sum is trivial exactly when all summands
are); then per summand the Alexander minor at $t=-1$ and at a generic field
element modulo $p=2^{61}-1$, the Kauffman bracket at a generic $A$ modulo $p$
(skipped when it exceeds 4096 frontier states), the exact Alexander polynomial
over $\mathbb Z[t]$, and finally the Khovanov scan with bit-packed morphisms,
compiled plans and min-fill pivots. Every filter is one-sided: only a mismatch
with the unknot's value gives a verdict, and that verdict is KNOTTED.

Not adopted: the exact integer bracket (equality with the unknot's value is
inconclusive whether it is computed modulo $p$ or over $\mathbb Z$, so it
cannot change a verdict, and it showed no measurable effect); replay commands
(useful, but not a speedup); tails longer than one crossing (slower, see below).

\subsection{Ablation}
\code{fast/ablation.py} measures each idea in isolation, again in fresh
processes (limit 300 s). First the scanner, in milliseconds:

\begin{center}\small
\IfFileExists{tables/ablation_scan.tex}{\begin{tabular}{@{}lrrrrr@{}}
\toprule
Scanner configuration & conway & braid3\_40 & braid4\_41 & braid6\_31 & braid5\_36 \\
\midrule
0.1 reference scanner & 42.9 & 1115 & 1264 & 612 & $>$300\,s \\
sets, LIFO, series inverse & 44.1 & 1044 & 1206 & 618 & $>$300\,s \\
sets, LIFO, self-inverse & 43.0 & 703 & 1113 & 618 & $>$300\,s \\
sets, min-fill & 43.0 & 594 & 845 & 553 & 10376 \\
bits, LIFO & 11.4 & 106 & 203 & 73.4 & 35622 \\
bits, min-fill (default) & 11.1 & 103 & 224 & 78.6 & 670 \\
bits, min-fill, tail 1 & 10.3 & 97.0 & 232 & 81.8 & 566 \\
bits, min-fill, tail 2 & 13.7 & 105 & 225 & 84.6 & 815 \\
\bottomrule
\end{tabular}
}{(table pending)}
\end{center}

Then the polynomial-time stages and the filters:

\begin{center}\small
\IfFileExists{tables/ablation_other.tex}{\begin{tabular}{@{}p{0.36\linewidth}p{0.23\linewidth}p{0.27\linewidth}r@{}}
\toprule
Task & new (ms) & old (ms) & ratio \\
\midrule
greedy order, 2048 crossings & heap O(n log n): 32.6 & 0.1 rescoring O(n$^2$): 26685 & 818$\times$ \\
descending test, 2048 crossings & linear sweep: 3.72 & 0.1 quadratic: 314 & 84$\times$ \\
R1/R2 reduction, 2048 crossings & incremental: 3.77 & 0.1 rebuild per move: 28089 & 7458$\times$ \\
T(3,61) rejection, 122 crossings & modular minor: 8.07 & 0.1 exact Z[t]: 5829 & 722$\times$ \\
recognize Conway & with Jones filter: 0.81 & without Jones filter: 10.0 & 12$\times$ \\
recognize Conway\#Conway\#Conway, no Jones & with factorization: 12.6 & without factorization: 6781 & 539$\times$ \\
rank of Conway\#Conway\#Conway & factored: 20.7 & one scan: 6972 & 337$\times$ \\
\bottomrule
\end{tabular}
}{(table pending)}
\end{center}

What the ablation shows:
\begin{itemize}
\item \textbf{The bit-packed algebra} is the across-the-board gain: 4 to 10
  times on every scan, with either pivot rule.
\item \textbf{Min-fill pivots} matter on exactly one kind of input, and there
  they are decisive. On the stress case the bit-packed scanner needs 35.6 s
  with LIFO pivots and 0.67 s with min-fill; the set-based scanner goes from
  more than 300 s to 10.4 s. On the other four inputs the pivot rule is
  neutral to slightly negative (the heap costs a few percent). The final
  complex is the same in all configurations, so the difference is fill-in
  during cancellation, not the size of the homology.
\item \textbf{Self-inverse units} save between nothing and 30\%, depending on
  how many cancellations involve a non-identity unit.
\item \textbf{The tail trick} is mixed: one crossing without cancellation is
  16\% faster on the stress case and 6\% faster on one other input, a few
  percent slower on the rest; two crossings are slower almost everywhere
  because the uncancelled complex quadruples. It is kept as an option
  (\code{tail}) with default zero.
\item \textbf{Heap ordering, incremental reduction, modular Alexander test}:
  three to four orders of magnitude at 2048 crossings and at $T(3,61)$. These
  are asymptotic improvements of polynomial stages: the old code spent 27 s on
  ordering and 28 s on R1 moves for a diagram that reduces to nothing.
\item \textbf{The linear descending test} needs an adversarial input to show
  (a 2048-crossing chain with one crossing flipped: 3.7 ms against 314 ms);
  on typical inputs the old test exits early and the two are equal.
\item \textbf{The Jones filter} turns recognition of the Conway knot from a
  10 ms Khovanov scan into a 0.8 ms bracket evaluation, and
  \textbf{factorization} turns a 6.8 s scan of three Conway summands into
  12.6 ms even with the Jones filter disabled.
\end{itemize}

\subsection{Complexity: what changed and what did not}
The polynomial stages improved asymptotically: scan ordering from $O(n^2)$ to
$O(n\log n)$ per start, the descending test from $O(n^2)$ to $O(n)$, R1/R2
reduction from one $O(n)$ revalidation per move to $O(1)$ per move, and the
Alexander rejection from fraction-free elimination over $\mathbb Z[t]$ to
$O(n^3)$ field operations.

Two statements with an exponential flavour are now \emph{proved} rather than
observed:
\begin{proposition}
The Jones filter costs $O\bigl(n\cdot\mathrm{Catalan}(w/2)\cdot\mathrm{poly}(w)\bigr)$
field operations for a scan of width $w$.
\end{proposition}
\begin{proof}
After each crossing the state is a map from crossingless matchings of at most
$w$ boundary points to $\mathbb F_p$; each entry spawns two successors.
\end{proof}
\begin{proposition}
If the visible factors of the reduced diagram have $n_1,\dots,n_k$ crossings,
recognition and the factored rank cost $\mathrm{poly}(n)+\sum_i2^{O(n_i)}$.
In particular the cost is polynomial when every $n_i=O(\log n)$ and
quasi-polynomial when every $n_i=O(\log^2n)$.
\end{proposition}
This is a statement about a restricted class of diagrams (several proposals
make it too), not about unknot recognition: a composite knot can be drawn
without any two-edge cut, and prime knots gain nothing. The family
$\#_k\,K11n34$, which Section \ref{sec:fast} used to show that a scan cannot be
subexponential if it must output the whole homology, is exactly the family
that factorization now handles in polynomial time; the obstruction moves to
diagrams in which the summands are hidden.

The general worst case of the exact fallback is unchanged, $2^{O(n)}$, and the
$\qp$ target of Section \ref{sec:missing} is no closer: nothing here touches
the hierarchy machinery.

\subsection{Verification}
The package now has 33 unit tests. The 14 new ones compare bit-packed with
set-based composition on random planar matchings, check every unit of a
three-circle endomorphism ring, run all scanner configurations against the 0.1
scanner, check that heap orders equal the 0.1 orders from every start, check
sign and filter invariance under half-turns of PD rows, check that no filter
fires on any unknot diagram and that every filter hit has reduced Khovanov
rank above one, check $33^k$ for $k=2,3,8$ Conway summands and equality of
factored and unfactored degree counts, replay simplifier traces, and
recompute the stress case (rank 2949, with and without the tail).

\section{A Rust implementation}
\label{sec:rust}

\subsection{Design}
\code{rust/} is a dependency-free crate (about 1900 lines) implementing the 0.2
pipeline except the exact Alexander polynomial over $\mathbb Z[t]$, which
would need big integers and on which exactness never depended. It was designed
for Rust rather than transliterated:
\begin{itemize}
\item matchings are interned, so objects, plans and caches hold 32-bit ids;
\item a morphism is a bitset over dot masks, a single inline \code{u64} when
  the Hom space has at most six circles (every scan of width at most twelve)
  and boxed words otherwise;
\item composition and transfer plans are a few integer masks per surface
  component, evaluated with popcounts;
\item memo tables use an Fx-style multiplicative hasher; with the standard
  SipHash, hashing dominated the profile;
\item objects live in flat vectors, the min-fill queue is a binary heap, and
  field arithmetic modulo $2^{61}-1$ uses 128-bit products;
\item factored ranks such as $33^{16}$ use a small base-$10^9$ integer.
\end{itemize}

\subsection{Verification}
Four unit tests (known ranks under both pivot rules, tails 0--2 and mirrors;
input validation; one-sidedness of the filters; the pipeline), and
\code{rust/cross\_check.py}, which compares ranks by cube degree in three
configurations and verdicts in two with the Python package on random braid
closures: 200 closures, no disagreement.

\subsection{Profile}
\code{rust/profile.py} records the binary's own timers (medians of five runs
inside one process, after parsing). The comparison columns of the table in
Section \ref{sec:proposals} show the Rust binary against the Python engines on
identical inputs. Larger random braid closures, raw scan, one run each with a
300-second budget:

\begin{center}\small
\IfFileExists{tables/rust_scaling.tex}{\begin{tabular}{@{}lrrrrrr@{}}
\toprule
Input & reduced rank & boundary & objects & total s & transfer s & eliminate s \\
\midrule
random 3-braid, 80 crossings & 5189 & 6 & 12240 & 0.235 & 0.101 & 0.122 \\
random 3-braid, 160 crossings & \multicolumn{6}{l}{budget of 300 s exhausted} \\
random 4-braid, 61 crossings & 643 & 8 & 3310 & 0.113 & 0.055 & 0.055 \\
random 4-braid, 81 crossings & \multicolumn{6}{l}{budget of 300 s exhausted} \\
random 5-braid, 48 crossings & 39 & 10 & 450 & 0.015 & 0.011 & 0.004 \\
random 5-braid, 60 crossings & 481 & 10 & 7398 & 4.663 & 0.525 & 4.108 \\
random 6-braid, 43 crossings & 177 & 10 & 530 & 0.023 & 0.014 & 0.008 \\
random 6-braid, 55 crossings & 1171 & 10 & 13294 & 1.484 & 0.434 & 1.021 \\
random 7-braid, 48 crossings & 29 & 10 & 1838 & 0.070 & 0.032 & 0.036 \\
random 8-braid, 49 crossings & 1709 & 12 & 3418 & 0.309 & 0.153 & 0.149 \\
\bottomrule
\end{tabular}
}{(table pending)}
\end{center}

\RUSTPROSE
